\documentclass[10pt,aspectratio=169]{beamer}
\usepackage{bma01}

\title[Aproximación lineal]{Aproximación lineal de una función en un punto}
\subtitle{Diferenciales, propagación de errores y método de Newton\\[2pt]\small Semana 9 -- Sesión 2}

\begin{document}

\begin{frame}
  \titlepage
\end{frame}

\begin{frame}{Objetivos de la sesión}
  \begin{enumerate}
    \item Construir la \resaltar{linealización} $L(x)=f(a)+f'(a)(x-a)$ y usarla para aproximar valores.
    \item Usar \resaltar{diferenciales} para estimar cambios y \resaltar{propagar errores} de medición.
    \item Obtener aproximaciones clásicas de la física ($\sen\theta\approx\theta$, $(1+x)^k\approx1+kx$).
    \item Aplicar la linealización para resolver ecuaciones: \resaltar{método de Newton}.
  \end{enumerate}
  \vspace{1mm}
  {\small\textbf{Competencia 3 del sílabo:} Aplica la derivada y su utilidad en problemas de ciencia e ingeniería.}
  \begin{block}{Plan (2 horas; incluye verificación previa de 5 min)}
    \begin{tabular}{ll}
      0:00 -- 0:30 & Linealización: idea geométrica y ejemplos numéricos\\
      0:30 -- 1:00 & Diferenciales y propagación de errores\\
      1:00 -- 1:20 & Aproximaciones en física\\
      1:20 -- 1:45 & Método de Newton\\
      1:45 -- 2:00 & Actividad colaborativa, resumen y cierre metacognitivo
    \end{tabular}
  \end{block}
\end{frame}

\begin{frame}{Aula invertida: preparación y verificación}
  \small
  \begin{previa}[antes de la sesión]
    \begin{enumerate}
      \item Lea sobre linealización, diferenciales y el método de Newton.
      \item Calcule la linealización de $f(x)=\sqrt x$ en $a=25$.
      \item Pregunta guía: ¿cuándo es confiable la aproximación lineal?
    \end{enumerate}
  \end{previa}
  \begin{block}{Verificación al inicio (5 min, en pares)}
    \begin{enumerate}
      \item Escriba la linealización $L(x)$ de $f$ en $a$.
      \item ¿Cuánto es $df$ si $f(x)=x^2$?
    \end{enumerate}
    Los pares discuten y votan; el docente aclara solo lo que genera dudas y ajusta el ritmo de la sesión.
  \end{block}
\end{frame}

\section{Linealización}

\begin{frame}{Una curva derivable es ``casi recta'' de cerca}
  \centering
  \begin{tikzpicture}
    \begin{axis}[bma, name=a1, xmin=0, xmax=8, ymin=0, ymax=3.2, width=5.3cm, height=4.4cm, xtick={4}, ytick={2}, samples=100, title={\scriptsize $[0,8]$}]
      \addplot[curva, domain=0:8] {sqrt(x)};
      \addplot[curva2, domain=0:8] {2+(x-4)/4};
    \end{axis}
    \begin{axis}[bma, at={(a1.east)}, anchor=west, xshift=0.5cm, name=a2, xmin=3, xmax=5, ymin=1.7, ymax=2.3, width=5.3cm, height=4.4cm, xtick={4}, ytick={2}, samples=100, title={\scriptsize $[3,5]$}]
      \addplot[curva, domain=3:5] {sqrt(x)};
      \addplot[curva2, domain=3:5] {2+(x-4)/4};
    \end{axis}
    \begin{axis}[bma, at={(a2.east)}, anchor=west, xshift=0.5cm, xmin=3.9, xmax=4.1, ymin=1.97, ymax=2.03, width=5.3cm, height=4.4cm, xtick={4}, ytick={2}, samples=100, title={\scriptsize $[3.9,4.1]$}, scaled ticks=false]
      \addplot[curva, domain=3.9:4.1] {sqrt(x)};
      \addplot[curva2, domain=3.9:4.1] {2+(x-4)/4};
    \end{axis}
  \end{tikzpicture}

  \vspace{2mm}
  \begin{idea}
    Al hacer ``zoom'' en $(4,2)$, la gráfica de $\sqrt x$ se vuelve indistinguible de su recta tangente. La tangente es la
    \textbf{mejor aproximación lineal} de $f$ cerca de $a$.
  \end{idea}
\end{frame}

\begin{frame}{Linealización}
  \small
  \begin{definicion}
    La \resaltar{linealización} de $f$ en $a$ es
    \[ L(x)=f(a)+f'(a)(x-a),\qquad f(x)\approx L(x)\ \text{para } x\text{ cerca de } a. \]
  \end{definicion}
  \begin{teorema}[es la mejor aproximación lineal]
    Si $f$ es derivable en $a$: \ $f(x)-L(x)=\varepsilon(x)\,(x-a)$ con $\varepsilon(x)\to0$ cuando $x\to a$.
    \ Ninguna otra recta por $(a,f(a))$ tiene esta propiedad.
  \end{teorema}
  \begin{ejemplo}[1]
    $f(x)=\sqrt x$, $a=4$: \ $L(x)=2+\frac14(x-4)$.
    \begin{center}
      \begin{tabular}{c|c|c|c}
        $x$ & $L(x)$ & $\sqrt x$ & error\\ \hline
        $4.1$ & $2.025$ & $2.024846$ & $1.5\times10^{-4}$\\
        $3.9$ & $1.975$ & $1.974842$ & $1.6\times10^{-4}$\\
        $5$ & $2.25$ & $2.236068$ & $1.4\times10^{-2}$
      \end{tabular}
    \end{center}
    El error crece al alejarse de $a$ (aproximadamente como $(x-a)^2$).
  \end{ejemplo}
\end{frame}

\begin{frame}{Aproximaciones lineales notables (cerca de $0$)}
  \begin{columns}
    \begin{column}{0.5\textwidth}
      \renewcommand{\arraystretch}{1.5}
      \begin{tabular}{l|l}
        $f(x)$ & $L(x)$ en $a=0$\\ \hline
        $\sen x$ & $x$\\
        $\tan x$ & $x$\\
        $\cos x$ & $1$\\
        $e^x$ & $1+x$\\
        $\ln(1+x)$ & $x$\\
        $(1+x)^k$ & $1+kx$
      \end{tabular}
    \end{column}
    \begin{column}{0.5\textwidth}
      \begin{ejemplo}[2]
        \begin{itemize}
          \item $\sen(0.1)\approx0.1$ \ (real $0.099833$).
          \item $e^{0.1}\approx1.1$ \ (real $1.105171$).
          \item $\sqrt{1.02}=(1+0.02)^{1/2}\approx1.01$ \ (real $1.009950$).
          \item $\sqrt[3]{26}=3\left(1-\frac1{27}\right)^{1/3}\approx3\left(1-\frac1{81}\right)\approx2.963$ \ (real $2.962496$).
        \end{itemize}
      \end{ejemplo}
    \end{column}
  \end{columns}
  \begin{observacion}
    La linealización de $\cos x$ en $0$ es la constante $1$: para una mejor aproximación se necesitan términos cuadráticos
    ($1-\frac{x^2}{2}$): \textbf{polinomio de Taylor} (Semana 12).
  \end{observacion}
\end{frame}

\section{Diferenciales}

\begin{frame}{Diferenciales}
  \begin{columns}
    \begin{column}{0.55\textwidth}
      \begin{definicion}
        Si $y=f(x)$, el \resaltar{diferencial} de $x$ es $dx$ (un incremento cualquiera) y el de $y$ es
        \[ dy=f'(x)\,dx. \]
      \end{definicion}
      \begin{itemize}
        \item $\Delta y=f(x+dx)-f(x)$: cambio \textbf{real}.
        \item $dy$: cambio a lo largo de la \textbf{tangente}.
        \item $\Delta y\approx dy$ si $dx$ es pequeño.
      \end{itemize}
      \begin{ejemplo}[3]
        $y=x^3$, $x=2$, $dx=0.05$: \ $dy=3\cdot4\cdot0.05=0.6$, \ $\Delta y=2.05^3-8=0.615125$.
      \end{ejemplo}
    \end{column}
    \begin{column}{0.45\textwidth}
      \centering
      \begin{tikzpicture}
        \begin{axis}[bma, xmin=0, xmax=3.3, ymin=0, ymax=4.5, width=6.3cm, height=5.6cm, xtick={1.2,2.4}, xticklabels={$x$,$x+dx$}, ytick=\empty, samples=100]
          \addplot[curva, domain=0.2:3] {0.45*x^2+0.5};
          \addplot[curva2, domain=0.5:3] {0.45*1.44+0.5+1.08*(x-1.2)};
          \draw[dashed] (axis cs:1.2,1.148) -- (axis cs:2.4,1.148);
          \draw[uniVerde, very thick] (axis cs:2.4,1.148) -- (axis cs:2.4,2.444);
          \draw[uniAzul, very thick] (axis cs:2.48,1.148) -- (axis cs:2.48,3.092);
          \node[uniVerde] at (axis cs:2.15,1.8) {\scriptsize $dy$};
          \node[uniAzul] at (axis cs:2.75,2.2) {\scriptsize $\Delta y$};
          \node at (axis cs:1.8,0.95) {\scriptsize $dx$};
        \end{axis}
      \end{tikzpicture}
    \end{column}
  \end{columns}
\end{frame}

\begin{frame}{Propagación de errores de medición}
  \small
  \begin{ejemplo}[4]
    El radio de una esfera se mide como $r=10$ cm con error máximo $\pm0.05$ cm. Estime el error en el volumen calculado.
    \[ V=\frac43\pi r^3\ \Rightarrow\ dV=4\pi r^2\,dr=4\pi(100)(0.05)\approx62.8\ \text{cm}^3. \]
    Error \textbf{relativo}: \ $\dfrac{dV}{V}=3\,\dfrac{dr}{r}=3(0.005)=0.015$ \ (1.5 por ciento).
  \end{ejemplo}
  \begin{proposicion}[regla práctica]
    Si $y=Cx^n$, entonces $\dfrac{dy}{y}=n\,\dfrac{dx}{x}$: \ el error relativo se multiplica por el exponente.
  \end{proposicion}
  \begin{ejemplo}[5: laboratorio]
    Con un péndulo, $g=\dfrac{4\pi^2L}{T^2}$. Si $T$ se mide con error relativo $1$ por ciento, el error relativo en $g$ debido a $T$ es
    $\approx2$ por ciento. ¡Hay que medir el periodo con más cuidado que la longitud!
  \end{ejemplo}
\end{frame}

\section{Aproximaciones en física}

\begin{frame}{Aproximaciones lineales en física}
  \begin{ejemplo}[6: péndulo]
    La ecuación exacta $\theta''=-\frac gL\sen\theta$ se ``linealiza'' con $\sen\theta\approx\theta$: \ $\theta''=-\frac gL\theta$ (oscilador armónico).
    Para $\theta=10^\circ=0.1745$ rad, el error relativo de $\sen\theta\approx\theta$ es solo $0.5$ por ciento.
  \end{ejemplo}
  \begin{ejemplo}[7: gravedad en altura]
    $g(h)=g_0\left(1+\frac hR\right)^{-2}\approx g_0\left(1-\frac{2h}{R}\right)$. \ A $h=10$ km ($R=6371$ km): \ $g\approx g_0(1-0.0031)$.
  \end{ejemplo}
  \begin{ejemplo}[8: relatividad]
    $E_c=mc^2\left[\left(1-\frac{v^2}{c^2}\right)^{-1/2}-1\right]\approx mc^2\cdot\frac12\frac{v^2}{c^2}=\frac12mv^2$. \
    ¡La mecánica clásica es la linealización de la relativista!
  \end{ejemplo}
\end{frame}

\section{Método de Newton}

\begin{frame}{Método de Newton--Raphson}
  \begin{columns}
    \begin{column}{0.55\textwidth}
      Para resolver $f(x)=0$: reemplazar $f$ por su linealización en $x_n$ y resolver $L(x)=0$:
      \[ \boxed{x_{n+1}=x_n-\frac{f(x_n)}{f'(x_n)}} \]
      \begin{ejemplo}[9]
        $f(x)=x^3-x-1$ (Semana 5, bisección), $x_0=1$:
        \begin{center}
          \small
          \begin{tabular}{c|l}
            $n$ & $x_n$\\ \hline
            0 & 1\\
            1 & 1.5\\
            2 & 1.347826\\
            3 & 1.325200\\
            4 & 1.324718
          \end{tabular}
        \end{center}
        ¡6 cifras correctas en 4 pasos! (Bisección necesitaría unos 20.)
      \end{ejemplo}
    \end{column}
    \begin{column}{0.45\textwidth}
      \centering
      \begin{tikzpicture}
        \begin{axis}[bma, xmin=0.8, xmax=1.7, ymin=-1.3, ymax=2, width=6.3cm, height=5.5cm, xtick={1,1.5}, xticklabels={$x_0$,$x_1$}, ytick=\empty, samples=100, domain=0.85:1.65]
          \addplot[curva] {x^3-x-1};
          \addplot[curva2, domain=0.9:1.6] {-1+2*(x-1)};
          \addplot[uniVerde, thick, domain=1.25:1.65] {0.875+5.75*(x-1.5)};
          \draw[dashed] (axis cs:1,0) -- (axis cs:1,-1);
          \draw[dashed] (axis cs:1.5,0) -- (axis cs:1.5,0.875);
          \addplot[only marks, mark=*, black] coordinates {(1,-1) (1.5,0.875) (1.3478,0)};
        \end{axis}
      \end{tikzpicture}

      \small Cada paso sigue la tangente hasta el eje $x$.
    \end{column}
  \end{columns}
\end{frame}

\begin{frame}{Newton: más ejemplos y cuidados}
  \begin{ejemplo}[10: raíces cuadradas (Herón)]
    Para $f(x)=x^2-a$: \ $x_{n+1}=x_n-\dfrac{x_n^2-a}{2x_n}=\dfrac12\left(x_n+\dfrac a{x_n}\right)$.
    \ Con $a=2$, $x_0=1$: \ $1.5$, $1.416667$, $1.414216$, \dots
  \end{ejemplo}
  \begin{observacion}[¿cuándo falla?]
    \begin{itemize}
      \item Si $f'(x_n)\approx0$, la tangente es casi horizontal y el siguiente punto se va lejos.
      \item Con mala elección de $x_0$ puede oscilar o converger a otra raíz. Ej.: $f(x)=\sqrt[3]x$ diverge desde cualquier $x_0\ne0$.
    \end{itemize}
  \end{observacion}
  \begin{idea}
    Bisección es lenta pero segura (Bolzano); Newton es rapidísimo cerca de la raíz pero necesita un buen punto de partida.
  \end{idea}
\end{frame}

\begin{frame}{Actividad colaborativa: Estimación sin calculadora}
  \small
  \begin{activa}[Trabajo colaborativo]
    Estime con linealización $\sqrt{26}$, $\sqrt[3]{8.1}$, $\ln 1.05$ y $\sen 31^\circ$. Luego calcule el \textbf{error relativo} con una calculadora y explique por qué unas aproximaciones son mejores. Además: el radio de una esfera mide $10\pm0.1$ cm; estime el error en el volumen con diferenciales.
  \end{activa}
  \begin{itemize}
    \item \textbf{Organización:} equipos de 3 a 4 con roles rotativos (coordinador, redactor, verificador, portavoz); 12--15 minutos.
    \item \textbf{Producto:} Tabla de estimaciones, errores y propagación de error.
    \item \textbf{Retroalimentación:} el docente recorre los equipos, hace preguntas y no da la solución hasta la puesta en común.
  \end{itemize}
\end{frame}

\section{Ejercicios}

\begin{frame}{Ejercicios propuestos}
  \small
  \begin{ejercicio}
    \begin{enumerate}
      \item Halle la linealización y aproxime: (a) $\sqrt[4]{16.2}$, \ (b) $\ln(0.98)$, \ (c) $\cos 61^\circ$, \ (d) $\frac{1}{0.997}$, \ (e) $e^{-0.05}$.
      \item Pruebe que $\sqrt{1+x}\approx1+\frac x2$ y diga para qué $x$ el error es menor que $0.01$ (use una gráfica o tabla).
      \item El lado de un cubo mide $30$ cm con error de $\pm0.1$ cm. Estime el error máximo en su volumen y en su área superficial, absolutos y relativos.
      \item La resistencia de un alambre es $R=\frac{k}{r^2}$. Si el radio se mide con error relativo de $2$ por ciento, ¿cuál es el error relativo en $R$?
      \item Use Newton para aproximar $\sqrt[3]{10}$ y la raíz de $\cos x=x$ con 5 cifras decimales.
      \item Un satélite está a $h$ km de altura. Use la aproximación lineal para estimar la distancia al horizonte $d=\sqrt{2Rh+h^2}$ cuando $h\ll R$.
      \item Muestre que la ley de los gases ideales da $\frac{dV}{V}=\frac{dT}{T}-\frac{dP}{P}$. Interprete.
    \end{enumerate}
  \end{ejercicio}
\end{frame}

\begin{frame}{Metacognición y uso responsable de la IA}
  \small
  \begin{columns}[t]
    \begin{column}{0.5\textwidth}
      \begin{metacog}[cierre, 3 min]
        \begin{enumerate}
          \metacomun
          \item ¿Cómo decido qué punto $a$ elegir para linealizar? ¿Qué criterio uso?
        \end{enumerate}
      \end{metacog}
    \end{column}
    \begin{column}{0.5\textwidth}
      \begin{iabox}
        Pida a una IA estimar $\sen 31^\circ$ por linealización. Revise si usó radianes y el punto base correcto.\\[2pt]
        \textit{Regla:} primero intente usted; luego use la herramienta para contrastar, nunca para sustituir su razonamiento. Anote el cambio en su bitácora.
      \end{iabox}
    \end{column}
  \end{columns}
\end{frame}

\begin{frame}{Resumen}
  \begin{itemize}
    \item $L(x)=f(a)+f'(a)(x-a)$: la recta tangente aproxima $f$ cerca de $a$.
    \item Notables: $\sen x\approx x$, $e^x\approx1+x$, $\ln(1+x)\approx x$, $(1+x)^k\approx1+kx$.
    \item Diferencial $dy=f'(x)dx\approx\Delta y$. Error relativo: $\frac{dy}{y}=n\frac{dx}{x}$ si $y=Cx^n$.
    \item Newton: $x_{n+1}=x_n-\frac{f(x_n)}{f'(x_n)}$; convergencia muy rápida si $x_0$ es bueno.
  \end{itemize}
  \vspace{3mm}
  \textbf{Próxima sesión:} derivadas de funciones inversas (trigonométricas e hiperbólicas inversas).
\end{frame}


\begin{frame}{Evaluación y trabajo posterior}
  \small
  \begin{evalbox}[Práctica del miércoles de la semana 9: PD5]
    \textbf{Cobertura:} semanas 7 y 9. \textbf{Naturaleza:} práctica dirigida en equipo, con retroalimentación inmediata.
  \end{evalbox}
  \begin{itemize}
    \item \criterios
    \item \declaracionia
    \item \textbf{Portafolio:} guarde el producto de la actividad, la bitácora de metacognición y los ejercicios propuestos de esta sesión; son evidencia de progreso y base de la retroalimentación.
  \end{itemize}
\end{frame}

\end{document}
